\section{Implementation}

\subsection{Hyperparameters}

During sampling, we perform DDIM sampling \citep{ddim} with $50$ denoising steps, setting DDIM eta to 0. The inference resolution is fixed at $320\times512$ pixels and the video length is $16$ frames in the normal setting. The video length of longer inference is $64$ frames and the inference resolution of larger inference is $640\times512$ pixels. The scale of the classifier-free guidance is set to $12$. $\alpha$ in Equation~\ref{eq:ca} is $\frac{0.25}{len\_target\_prompts \times proportion\_target\_box}$ and $\beta$ in Equation~\ref{eq:sata} is $0.01$.
%

For quantitative comparison, we generate a total of $896$ videos for each inference method, utilizing $56$ prompts. We initialize $16$ random initial noises for each prompt for direct inference. For trajectory control methods, each prompt is applied to $8$ different trajectories with $2$ random initial noises. 

In the user study, we mixed our generated videos with those generated by the other three baselines. A total of $27$ users were asked to pick the best one according to the trajectory alignment, video-text alignment, and video quality, respectively.



\subsection{Prompts}

Our prompts are mostly extended from previous baselines\citep{jain2023peekaboo, ma2023trailblazer} but replace some prompts that conflict with object movement, like standing or lying.

\begin{itemize}
    \item \textbf{A woodpecker} climbing up a tree trunk.
    \item \textbf{A squirrel} descending a tree after gathering nuts.
    \item \textbf{A bird} diving towards the water to catch fish.
    \item \textbf{A frog} leaping up to catch a fly.
    \item \textbf{A parrot} flying upwards towards the treetops.
    \item \textbf{A squirrel} jumping from one tree to another.
    \item \textbf{A rabbit} burrowing downwards into its warren. 
    \item \textbf{A satellite} orbiting Earth in outer space.
    \item \textbf{A skateboarder} performing tricks at a skate park. 
    \item \textbf{A leaf} falling gently from a tree.
    \item \textbf{A paper plane} gliding in the air.
    \item \textbf{A bear} climbing down a tree after spotting a threat. 
    \item \textbf{A duck} diving underwater in search of food.
    \item \textbf{A kangaroo} hopping down a gentle slope.
    \item \textbf{An owl} swooping down on its prey during the night.
    \item \textbf{A hot air balloon} drifting across a clear sky.
    \item \textbf{A red double-decker bus} moving through London streets. 
    \item \textbf{A jet plane} flying high in the sky.
    \item \textbf{A helicopter} hovering above a cityscape.
    \item \textbf{A roller coaster} looping in an amusement park.
    \item \textbf{A streetcar} trundling down tracks in a historic district.
    \item \textbf{A rocket} launching into space from a launchpad.
    \item \textbf{A deer} walking in a snowy field.
    \item \textbf{A horse} grazing in a meadow.
    \item \textbf{A fox} running in a forest clearing.
    \item \textbf{A swan} floating gracefully on a lake.
    \item \textbf{A panda} walking and munching bamboo in a bamboo forest.
    \item \textbf{A penguin} walking on an iceberg.
    \item \textbf{A lion} walking in the savanna grass.
    \item \textbf{An owl} flying in a tree at night.
    \item \textbf{A dolphin} just breaking the ocean surface.
    \item \textbf{A camel} walking in a desert landscape.
    \item \textbf{A kangaroo} jumping in the Australian outback.
    \item \textbf{A colorful hot air balloon} tethered to the ground.
    \item \textbf{A corgi} running on the grassland on the grassland.
    \item \textbf{A corgi} running on the grassland in the snow.
    \item \textbf{A man} in gray clothes running in the summer.
    \item \textbf{A knight} riding a horse on a race course.
    \item \textbf{A horse} galloping on a street.
    \item \textbf{A lion} running on the grasslands.
    \item \textbf{A dog} running across the garden, photorealistic, 4k.
    \item \textbf{A tiger} walking in the forest, photorealistic, 4k, high definition.
    \item \textbf{Iron Man} surfing on the sea.
    \item \textbf{A tiger} running in the forest, photorealistic, 4k, high definition.
    \item \textbf{A horse} running, photorealistic, 4k, volumetric lighting unreal engine.
    \item \textbf{A panda} surfing in the universe.
    \item \textbf{A chihuahua} in an astronaut suit floating in the universe, cinematic lighting, glow effect.
    \item \textbf{An astronaut} waving his hands on the moon.
    \item \textbf{A horse} galloping through a meadow.
    \item \textbf{A bear} running in the ruins, photorealistic, 4k, high definition.
    \item \textbf{A barrel} floating in a river.
    \item \textbf{A dark knight} riding a horse on the grassland.
    \item \textbf{A wooden boat} moving on the sea.
    \item \textbf{A red car} turning around on a countryside road, photorealistic, 4k.
    \item \textbf{A majestic eagle} soaring high above the treetops, surveying its territory.
    \item \textbf{A bald eagle} flying in the blue sky.
    
\end{itemize}

\section{Longer and Larger Video Generation}

\begin{figure*}
\centering
\includegraphics[width=0.98\linewidth]{Figures/fig_long.pdf}
\caption{\textbf{Longer video generation.} Longer video generation allows us to plan some complex trajectories. FreeTraj succeeds in generating rich motion trajectories in long videos.}
\label{fig:long}
\end{figure*}

\subsection{Related Work of Long Video Generation.}
%
Long video generation is a challenging but important problem in video generation.
TATs~\citep{ge2022long}, longvideoGAN~\citep{ge2022long}, LVDM~\citep{lvdm}, and flexible diffusion~\citep{harvey2022flexible} achieve long video generation in small domains and without textual guidance.
Phenaki~\citep{villegas2022phenaki}, NUWA-Infinity~\citep{liang2022nuwa}, NUWA-XL~\citep{yin2023nuwa}, and Sora~\citep{sora} are text-guided long video generation approaches for open-domain generation.
Animate-A-Story~\citep{he2023animate} achieves multi-scene long video generation via character consistency.
Streamingt2v~\citep{streamingt2v} and FlexiFilm~\citep{flexifilm} are training-based methods that train a conditional module on top of pre-trained video diffusion models conditioning on previous-frames.
Genlvideo~\citep{wang2023gen} and FreeNoise~\citep{freenoise} are recently proposed tuning-free methods for generating longer videos based on pre-trained video diffusion models to extend their generated length.
In this work, we propose a tuning-free approach for long video generation based on long-term trajectory control.

\subsection{Results of Longer Generation}

FreeTraj can be integrated into the longer video generation framework FreeNoise~\citep{freenoise}. With the help of some technical points proposed by FreeNoise, our FreeTraj successfully generated trajectory-controllable long videos. As shown in Figure~\ref{fig:long}, we plan two complex paths and FreeTraj succeeds in generating rich motion trajectories in long videos.

\subsection{Results of Larger Generation}

When we directly use pre-trained video diffusion models to generate videos with higher resolutions compared to those in training, they will easily generate results with duplicated main objects anywhere~\cite {he2023scalecrafter}. However, FreeTraj will plan the trajectory for the main object, and information about the main object will be reduced out of the target areas. Therefore, the duplication phenomenon will be suppressed by FreeTraj (Figure~\ref{fig:large}).

\begin{figure*}
\centering
\includegraphics[width=0.98\linewidth]{Figures/fig_large.pdf}
\caption{\textbf{Larger video generation.} Directly generating larger videos will easily lead to the results with duplicated main objects anywhere. FreeTraj plans the trajectory for the main object and suppresses the duplication phenomenon.}
\label{fig:large}
\end{figure*}

\section{More Observations}

\subsection{More About Noise Flow}

Here we show another direction of noise flow. Instead of randomly sampling initial noises for all frames, we only sample the noise for the first frame. Then we move the noise from the bottom-left to the top-right with stride $2$ and repeat this operation until we get initial noises $z_{T}^{flow}$ for all frames. Specially, initial noise $\epsilon$ for each frame $f$ in position $[i, j]$ is:
%
\begin{equation}
\epsilon[i, j]^{f} = \epsilon[(i+2) \ (\mathrm{mod} \ H), (j-2) \ (\mathrm{mod} \ W)]^{f-1}.
\end{equation}
%
After denoising $z_{T}^{flow}$, results in Figure~\ref{fig:moveanalysis_re} show that objects and textures in the video also flow in the same direction (bottom-left to top-right). This phenomenon verifies that the trajectory of the initial noises can have an impact on the motion trajectory of the generated result. When the proportion of high-frequency noise resampled increases, the visual quality is significantly improved. Correspondingly, the flow phenomenon is weakened.

\begin{figure*} %[t]
% \vspace{-0.1in}
\centering
\includegraphics[width=0.95\linewidth]{Figures/fig_moveanalysis_re.pdf}
\vspace{-1.0em}
\caption{\textbf{Noise resampling of initial high-frequency components.} Gradually increasing the proportion of resampled high-frequency information in the frame-wise shared noises can significantly reduce the artifact in the generated video. However, this also leads to a gradual loss in trajectory control ability. A resampling percentage of 75\% strikes a better balance between maintaining control and improving the quality of the generated video.}
\vspace{-1.0em}
\label{fig:moveanalysis_re}
\end{figure*}

\subsection{Attention Isolation in Temporal Dimension}

Usually, we initialize $16$ frames of random noises independently. Instead of normal sampling, we try partial repeated sampling by partially repeating some initial noises:
%
\begin{equation}
\begin{aligned}
&\text{Normal Sampling: } [\epsilon_1, \epsilon_2, \epsilon_3, \epsilon_4, \epsilon_5, \epsilon_6, \epsilon_7, \epsilon_8, \epsilon_9, \epsilon_{10}, \epsilon_{11}, \epsilon_{12}, \epsilon_{13}, \epsilon_{14}, \epsilon_{15}, \epsilon_{16}], \\
&\text{Partial Repeated Sampling: } [\mathbf{\epsilon_1}, \mathbf{\epsilon_1}, \mathbf{\epsilon_1}, \mathbf{\epsilon_1}, \epsilon_2, \epsilon_3, \epsilon_4, \epsilon_5, \epsilon_6, \epsilon_{7}, \epsilon_{8}, \epsilon_{9}, \mathbf{\epsilon_{10}}, \mathbf{\epsilon_{10}}, \mathbf{\epsilon_{10}}, \mathbf{\epsilon_{10}}].
\end{aligned}
\end{equation}

Since spatio-temporal correlations in the low-frequency components of initial noises will guide the trajectory of generated objects, partial repeated sampling for initial noises will bring typical motion mode. As shown in Figure~\ref{fig:tempnoise} (b), the owl is stationary in the beginning and ending frames and only has significant action in the middle frames. However, due to the attention isolation, frames of generated results have obvious artifacts. We visualize one heat map of temporal attention and find that stationary frames mainly pay attention to frames with the same initial noises. When calculating the attention weights received by isolated frames, manually splitting a portion of attention weights from isolated frames to other frames will remove artifacts. As shown in Figure~\ref{fig:tempnoise} (c), an owl is well generated and its motion still fits the mode in (b).

\begin{figure*} %[t]
% \vspace{-0.1in}
\centering
\includegraphics[width=0.95\linewidth]{Figures/fig_tempnoise.pdf}
\vspace{-1.0em}
\caption{\textbf{Attention isolation in temporal dimension.} Compared to normal sampling for initial noises (a), partial repeated sampling will lead to significant attention isolation in the temporal dimension and bring strong artifacts (b). When calculating the attention weights received by isolated frames, manually splitting a portion of attention weights from isolated frames to other frames will remove artifacts (c).}
\vspace{-1.0em}
\label{fig:tempnoise}
\end{figure*}
